% paper25-metaobject-analysis.tex — Judgment Geometry: Metaobject Protocol Analysis
% Paper 25 of the Judgment Geometry series.
% Compile: pdflatex -interaction=nonstopmode paper25-metaobject-analysis.tex
\documentclass[11pt]{article}
% jugeo-common.tex — shared preamble for all Judgment Geometry papers
% Include via: % jugeo-common.tex — shared preamble for all Judgment Geometry papers
% Include via: % jugeo-common.tex — shared preamble for all Judgment Geometry papers
% Include via: \input{jugeo-common}

% ── Packages ────────────────────────────────────────────────────────────
\usepackage{amsmath,amssymb,amsthm,mathtools}
\usepackage{tikz-cd}
\usepackage{listings}
\usepackage{xcolor}
\usepackage{xspace}
\usepackage{booktabs}
\usepackage{multirow}
\usepackage{enumitem}
\usepackage[expansion=false]{microtype}
\usepackage{hyperref}
\usepackage{cleveref}
% \usepackage{thmtools}  % disabled: not available in this TeX installation
\usepackage{float}
\usepackage{caption}
\usepackage{subcaption}
\usepackage{tabularx}
\usepackage{stmaryrd}       % \llbracket, \rrbracket
\usepackage{bbm}            % \mathbbm{1}

% ── Page geometry ───────────────────────────────────────────────────────
\usepackage[margin=1in]{geometry}

% ── Theorem environments ────────────────────────────────────────────────
\theoremstyle{plain}
\newtheorem{theorem}{Theorem}[section]
\newtheorem{lemma}[theorem]{Lemma}
\newtheorem{proposition}[theorem]{Proposition}
\newtheorem{corollary}[theorem]{Corollary}

\theoremstyle{definition}
\newtheorem{definition}[theorem]{Definition}
\newtheorem{example}[theorem]{Example}
\newtheorem{construction}[theorem]{Construction}

\theoremstyle{remark}
\newtheorem{remark}[theorem]{Remark}
\newtheorem{notation}[theorem]{Notation}

% ── Python listings style ──────────────────────────────────────────────
\definecolor{jugeoBlue}{HTML}{2563EB}
\definecolor{jugeoGreen}{HTML}{059669}
\definecolor{jugeoRed}{HTML}{DC2626}
\definecolor{jugeoPurple}{HTML}{7C3AED}
\definecolor{jugeoGray}{HTML}{6B7280}
\definecolor{jugeoBg}{HTML}{F8FAFC}

\lstdefinestyle{jugeo-python}{
  language=Python,
  basicstyle=\ttfamily\small,
  keywordstyle=\color{jugeoBlue}\bfseries,
  stringstyle=\color{jugeoGreen},
  commentstyle=\color{jugeoGray}\itshape,
  numberstyle=\tiny\color{jugeoGray},
  backgroundcolor=\color{jugeoBg},
  numbers=left,
  numbersep=6pt,
  frame=single,
  framerule=0.4pt,
  rulecolor=\color{jugeoGray!40},
  breaklines=true,
  breakatwhitespace=true,
  showstringspaces=false,
  tabsize=4,
  xleftmargin=1.5em,
  framexleftmargin=1.5em,
  morekeywords={dataclass,frozen,field,Enum,IntEnum,auto,Any,
                Mapping,Sequence,Optional,match,case},
  literate={→}{{$\to$}}1 {←}{{$\leftarrow$}}1
           {≤}{{$\leq$}}1 {≥}{{$\geq$}}1 {≠}{{$\neq$}}1
           {∈}{{$\in$}}1 {∀}{{$\forall$}}1 {∃}{{$\exists$}}1
           {⊕}{{$\oplus$}}1 {⊖}{{$\ominus$}}1,
  escapeinside={(*@}{@*)},
}
\lstset{style=jugeo-python}

\lstdefinestyle{jugeo-output}{
  basicstyle=\ttfamily\small,
  backgroundcolor=\color{jugeoBg},
  frame=single,
  framerule=0.4pt,
  rulecolor=\color{jugeoGray!40},
  breaklines=true,
  showstringspaces=false,
  xleftmargin=1.5em,
  framexleftmargin=0.5em,
  numbers=none,
}

% ── JuGeo-specific macros ──────────────────────────────────────────────

% Judgment 8-tuple
\newcommand{\J}{\mathcal{J}}
\newcommand{\Jud}[8]{(#1,\,#2,\,#3,\,#4,\,#5,\,#6,\,#7,\,#8)}
\newcommand{\JudShort}[1]{\J_{#1}}

% Components
\newcommand{\coord}{\mathsf{c}}            % coordinate
\newcommand{\prop}{\varphi}                 % proposition
\newcommand{\carrier}{\mathcal{A}}          % carrier type
\newcommand{\evid}{\mathcal{E}}             % evidence bundle
\newcommand{\oblig}{\mathcal{O}}            % obligations
\newcommand{\obstr}{\mathcal{B}}            % obstructions
\newcommand{\trust}{\mathcal{T}}            % trust annotation
\newcommand{\prov}{\Pi}                     % provenance

% Trust algebra
\newcommand{\Talg}{\mathfrak{T}}            % trust ordered algebra
\newcommand{\Eadm}{\mathcal{E}_{\mathrm{adm}}}
\newcommand{\tleq}{\preceq}                 % trust ordering
\newcommand{\tjoin}{\oplus}                 % conservative join
\newcommand{\tatten}{\ominus}               % attenuation
\newcommand{\tpromote}{\uparrow_\pi}        % promotion
\newcommand{\tdemote}{\downarrow_\chi}       % demotion

% Trust tiers
\newcommand{\Tcontra}{\ensuremath{\mathsf{CONTRADICTED}}}
\newcommand{\Tunver}{\ensuremath{\mathsf{UNVERIFIED}}}
\newcommand{\Tcopilot}{\ensuremath{\mathsf{COPILOT}}}
\newcommand{\Toracle}{\ensuremath{\mathsf{ORACLE}}}
\newcommand{\Truntime}{\ensuremath{\mathsf{RUNTIME}}}
\newcommand{\Tsolver}{\ensuremath{\mathsf{SOLVER}}}
\newcommand{\Tproof}{\ensuremath{\mathsf{PROOF}}}

% Site and geometry
\newcommand{\Site}{\mathcal{S}}             % semantic site
\newcommand{\Cov}{\mathrm{Cov}}            % covering family
\newcommand{\Ob}{\mathrm{Ob}}              % objects
\newcommand{\Mor}{\mathrm{Mor}}            % morphisms
\newcommand{\res}{\mathrm{res}}            % restriction
\newcommand{\Cech}{\check{C}}              % Čech complex
\newcommand{\Hcoh}[1]{H^{#1}}             % cohomology group

% Descent
\newcommand{\descent}{\mathrm{desc}}
\newcommand{\glue}{\mathrm{glue}}
\newcommand{\overlap}[2]{#1 \cap #2}

% Semantic moves
\newcommand{\VERIFY}{\textsc{Verify}}
\newcommand{\CONSTRUCT}{\textsc{Construct}}
\newcommand{\NORMALIZE}{\textsc{Normalize}}
\newcommand{\GLUE}{\textsc{Glue}}
\newcommand{\RESTRICT}{\textsc{Restrict}}
\newcommand{\TRANSPORT}{\textsc{Transport}}
\newcommand{\REPAIR}{\textsc{Repair}}
\newcommand{\PROMOTE}{\textsc{Promote}}
\newcommand{\CHALLENGE}{\textsc{Challenge}}
\newcommand{\ESCALATE}{\textsc{Escalate}}

% Fragments
\newcommand{\QFLIA}{\texttt{QF\_LIA}}
\newcommand{\QFLRA}{\texttt{QF\_LRA}}
\newcommand{\QFBV}{\texttt{QF\_BV}}
\newcommand{\QFUF}{\texttt{QF\_UF}}

% Misc
\newcommand{\Python}{\textsc{Python}}
\newcommand{\JuGeo}{\textsc{JuGeo}}
\newcommand{\LEAN}{\textsc{Lean}\,4}
\newcommand{\Fstar}{F$^{\star}$}
\newcommand{\Coq}{\textsc{Coq}}
\newcommand{\Dafny}{\textsc{Dafny}}

% Common shortcuts
\newcommand{\ie}{i.e.\@\xspace}
\newcommand{\eg}{e.g.\@\xspace}
\newcommand{\cf}{cf.\@\xspace}
\newcommand{\wrt}{w.r.t.\@\xspace}

% Evaluation shortcuts
\newcommand{\numcases}{300}
\newcommand{\numspec}{100}
\newcommand{\numeq}{100}
\newcommand{\numbug}{100}
\newcommand{\medtime}{6.5\,ms}
\newcommand{\totaltime}{1.949\,s}


% ── Packages ────────────────────────────────────────────────────────────
\usepackage{amsmath,amssymb,amsthm,mathtools}
\usepackage{tikz-cd}
\usepackage{listings}
\usepackage{xcolor}
\usepackage{xspace}
\usepackage{booktabs}
\usepackage{multirow}
\usepackage{enumitem}
\usepackage[expansion=false]{microtype}
\usepackage{hyperref}
\usepackage{cleveref}
% \usepackage{thmtools}  % disabled: not available in this TeX installation
\usepackage{float}
\usepackage{caption}
\usepackage{subcaption}
\usepackage{tabularx}
\usepackage{stmaryrd}       % \llbracket, \rrbracket
\usepackage{bbm}            % \mathbbm{1}

% ── Page geometry ───────────────────────────────────────────────────────
\usepackage[margin=1in]{geometry}

% ── Theorem environments ────────────────────────────────────────────────
\theoremstyle{plain}
\newtheorem{theorem}{Theorem}[section]
\newtheorem{lemma}[theorem]{Lemma}
\newtheorem{proposition}[theorem]{Proposition}
\newtheorem{corollary}[theorem]{Corollary}

\theoremstyle{definition}
\newtheorem{definition}[theorem]{Definition}
\newtheorem{example}[theorem]{Example}
\newtheorem{construction}[theorem]{Construction}

\theoremstyle{remark}
\newtheorem{remark}[theorem]{Remark}
\newtheorem{notation}[theorem]{Notation}

% ── Python listings style ──────────────────────────────────────────────
\definecolor{jugeoBlue}{HTML}{2563EB}
\definecolor{jugeoGreen}{HTML}{059669}
\definecolor{jugeoRed}{HTML}{DC2626}
\definecolor{jugeoPurple}{HTML}{7C3AED}
\definecolor{jugeoGray}{HTML}{6B7280}
\definecolor{jugeoBg}{HTML}{F8FAFC}

\lstdefinestyle{jugeo-python}{
  language=Python,
  basicstyle=\ttfamily\small,
  keywordstyle=\color{jugeoBlue}\bfseries,
  stringstyle=\color{jugeoGreen},
  commentstyle=\color{jugeoGray}\itshape,
  numberstyle=\tiny\color{jugeoGray},
  backgroundcolor=\color{jugeoBg},
  numbers=left,
  numbersep=6pt,
  frame=single,
  framerule=0.4pt,
  rulecolor=\color{jugeoGray!40},
  breaklines=true,
  breakatwhitespace=true,
  showstringspaces=false,
  tabsize=4,
  xleftmargin=1.5em,
  framexleftmargin=1.5em,
  morekeywords={dataclass,frozen,field,Enum,IntEnum,auto,Any,
                Mapping,Sequence,Optional,match,case},
  literate={→}{{$\to$}}1 {←}{{$\leftarrow$}}1
           {≤}{{$\leq$}}1 {≥}{{$\geq$}}1 {≠}{{$\neq$}}1
           {∈}{{$\in$}}1 {∀}{{$\forall$}}1 {∃}{{$\exists$}}1
           {⊕}{{$\oplus$}}1 {⊖}{{$\ominus$}}1,
  escapeinside={(*@}{@*)},
}
\lstset{style=jugeo-python}

\lstdefinestyle{jugeo-output}{
  basicstyle=\ttfamily\small,
  backgroundcolor=\color{jugeoBg},
  frame=single,
  framerule=0.4pt,
  rulecolor=\color{jugeoGray!40},
  breaklines=true,
  showstringspaces=false,
  xleftmargin=1.5em,
  framexleftmargin=0.5em,
  numbers=none,
}

% ── JuGeo-specific macros ──────────────────────────────────────────────

% Judgment 8-tuple
\newcommand{\J}{\mathcal{J}}
\newcommand{\Jud}[8]{(#1,\,#2,\,#3,\,#4,\,#5,\,#6,\,#7,\,#8)}
\newcommand{\JudShort}[1]{\J_{#1}}

% Components
\newcommand{\coord}{\mathsf{c}}            % coordinate
\newcommand{\prop}{\varphi}                 % proposition
\newcommand{\carrier}{\mathcal{A}}          % carrier type
\newcommand{\evid}{\mathcal{E}}             % evidence bundle
\newcommand{\oblig}{\mathcal{O}}            % obligations
\newcommand{\obstr}{\mathcal{B}}            % obstructions
\newcommand{\trust}{\mathcal{T}}            % trust annotation
\newcommand{\prov}{\Pi}                     % provenance

% Trust algebra
\newcommand{\Talg}{\mathfrak{T}}            % trust ordered algebra
\newcommand{\Eadm}{\mathcal{E}_{\mathrm{adm}}}
\newcommand{\tleq}{\preceq}                 % trust ordering
\newcommand{\tjoin}{\oplus}                 % conservative join
\newcommand{\tatten}{\ominus}               % attenuation
\newcommand{\tpromote}{\uparrow_\pi}        % promotion
\newcommand{\tdemote}{\downarrow_\chi}       % demotion

% Trust tiers
\newcommand{\Tcontra}{\ensuremath{\mathsf{CONTRADICTED}}}
\newcommand{\Tunver}{\ensuremath{\mathsf{UNVERIFIED}}}
\newcommand{\Tcopilot}{\ensuremath{\mathsf{COPILOT}}}
\newcommand{\Toracle}{\ensuremath{\mathsf{ORACLE}}}
\newcommand{\Truntime}{\ensuremath{\mathsf{RUNTIME}}}
\newcommand{\Tsolver}{\ensuremath{\mathsf{SOLVER}}}
\newcommand{\Tproof}{\ensuremath{\mathsf{PROOF}}}

% Site and geometry
\newcommand{\Site}{\mathcal{S}}             % semantic site
\newcommand{\Cov}{\mathrm{Cov}}            % covering family
\newcommand{\Ob}{\mathrm{Ob}}              % objects
\newcommand{\Mor}{\mathrm{Mor}}            % morphisms
\newcommand{\res}{\mathrm{res}}            % restriction
\newcommand{\Cech}{\check{C}}              % Čech complex
\newcommand{\Hcoh}[1]{H^{#1}}             % cohomology group

% Descent
\newcommand{\descent}{\mathrm{desc}}
\newcommand{\glue}{\mathrm{glue}}
\newcommand{\overlap}[2]{#1 \cap #2}

% Semantic moves
\newcommand{\VERIFY}{\textsc{Verify}}
\newcommand{\CONSTRUCT}{\textsc{Construct}}
\newcommand{\NORMALIZE}{\textsc{Normalize}}
\newcommand{\GLUE}{\textsc{Glue}}
\newcommand{\RESTRICT}{\textsc{Restrict}}
\newcommand{\TRANSPORT}{\textsc{Transport}}
\newcommand{\REPAIR}{\textsc{Repair}}
\newcommand{\PROMOTE}{\textsc{Promote}}
\newcommand{\CHALLENGE}{\textsc{Challenge}}
\newcommand{\ESCALATE}{\textsc{Escalate}}

% Fragments
\newcommand{\QFLIA}{\texttt{QF\_LIA}}
\newcommand{\QFLRA}{\texttt{QF\_LRA}}
\newcommand{\QFBV}{\texttt{QF\_BV}}
\newcommand{\QFUF}{\texttt{QF\_UF}}

% Misc
\newcommand{\Python}{\textsc{Python}}
\newcommand{\JuGeo}{\textsc{JuGeo}}
\newcommand{\LEAN}{\textsc{Lean}\,4}
\newcommand{\Fstar}{F$^{\star}$}
\newcommand{\Coq}{\textsc{Coq}}
\newcommand{\Dafny}{\textsc{Dafny}}

% Common shortcuts
\newcommand{\ie}{i.e.\@\xspace}
\newcommand{\eg}{e.g.\@\xspace}
\newcommand{\cf}{cf.\@\xspace}
\newcommand{\wrt}{w.r.t.\@\xspace}

% Evaluation shortcuts
\newcommand{\numcases}{300}
\newcommand{\numspec}{100}
\newcommand{\numeq}{100}
\newcommand{\numbug}{100}
\newcommand{\medtime}{6.5\,ms}
\newcommand{\totaltime}{1.949\,s}


% ── Packages ────────────────────────────────────────────────────────────
\usepackage{amsmath,amssymb,amsthm,mathtools}
\usepackage{tikz-cd}
\usepackage{listings}
\usepackage{xcolor}
\usepackage{xspace}
\usepackage{booktabs}
\usepackage{multirow}
\usepackage{enumitem}
\usepackage[expansion=false]{microtype}
\usepackage{hyperref}
\usepackage{cleveref}
% \usepackage{thmtools}  % disabled: not available in this TeX installation
\usepackage{float}
\usepackage{caption}
\usepackage{subcaption}
\usepackage{tabularx}
\usepackage{stmaryrd}       % \llbracket, \rrbracket
\usepackage{bbm}            % \mathbbm{1}

% ── Page geometry ───────────────────────────────────────────────────────
\usepackage[margin=1in]{geometry}

% ── Theorem environments ────────────────────────────────────────────────
\theoremstyle{plain}
\newtheorem{theorem}{Theorem}[section]
\newtheorem{lemma}[theorem]{Lemma}
\newtheorem{proposition}[theorem]{Proposition}
\newtheorem{corollary}[theorem]{Corollary}

\theoremstyle{definition}
\newtheorem{definition}[theorem]{Definition}
\newtheorem{example}[theorem]{Example}
\newtheorem{construction}[theorem]{Construction}

\theoremstyle{remark}
\newtheorem{remark}[theorem]{Remark}
\newtheorem{notation}[theorem]{Notation}

% ── Python listings style ──────────────────────────────────────────────
\definecolor{jugeoBlue}{HTML}{2563EB}
\definecolor{jugeoGreen}{HTML}{059669}
\definecolor{jugeoRed}{HTML}{DC2626}
\definecolor{jugeoPurple}{HTML}{7C3AED}
\definecolor{jugeoGray}{HTML}{6B7280}
\definecolor{jugeoBg}{HTML}{F8FAFC}

\lstdefinestyle{jugeo-python}{
  language=Python,
  basicstyle=\ttfamily\small,
  keywordstyle=\color{jugeoBlue}\bfseries,
  stringstyle=\color{jugeoGreen},
  commentstyle=\color{jugeoGray}\itshape,
  numberstyle=\tiny\color{jugeoGray},
  backgroundcolor=\color{jugeoBg},
  numbers=left,
  numbersep=6pt,
  frame=single,
  framerule=0.4pt,
  rulecolor=\color{jugeoGray!40},
  breaklines=true,
  breakatwhitespace=true,
  showstringspaces=false,
  tabsize=4,
  xleftmargin=1.5em,
  framexleftmargin=1.5em,
  morekeywords={dataclass,frozen,field,Enum,IntEnum,auto,Any,
                Mapping,Sequence,Optional,match,case},
  literate={→}{{$\to$}}1 {←}{{$\leftarrow$}}1
           {≤}{{$\leq$}}1 {≥}{{$\geq$}}1 {≠}{{$\neq$}}1
           {∈}{{$\in$}}1 {∀}{{$\forall$}}1 {∃}{{$\exists$}}1
           {⊕}{{$\oplus$}}1 {⊖}{{$\ominus$}}1,
  escapeinside={(*@}{@*)},
}
\lstset{style=jugeo-python}

\lstdefinestyle{jugeo-output}{
  basicstyle=\ttfamily\small,
  backgroundcolor=\color{jugeoBg},
  frame=single,
  framerule=0.4pt,
  rulecolor=\color{jugeoGray!40},
  breaklines=true,
  showstringspaces=false,
  xleftmargin=1.5em,
  framexleftmargin=0.5em,
  numbers=none,
}

% ── JuGeo-specific macros ──────────────────────────────────────────────

% Judgment 8-tuple
\newcommand{\J}{\mathcal{J}}
\newcommand{\Jud}[8]{(#1,\,#2,\,#3,\,#4,\,#5,\,#6,\,#7,\,#8)}
\newcommand{\JudShort}[1]{\J_{#1}}

% Components
\newcommand{\coord}{\mathsf{c}}            % coordinate
\newcommand{\prop}{\varphi}                 % proposition
\newcommand{\carrier}{\mathcal{A}}          % carrier type
\newcommand{\evid}{\mathcal{E}}             % evidence bundle
\newcommand{\oblig}{\mathcal{O}}            % obligations
\newcommand{\obstr}{\mathcal{B}}            % obstructions
\newcommand{\trust}{\mathcal{T}}            % trust annotation
\newcommand{\prov}{\Pi}                     % provenance

% Trust algebra
\newcommand{\Talg}{\mathfrak{T}}            % trust ordered algebra
\newcommand{\Eadm}{\mathcal{E}_{\mathrm{adm}}}
\newcommand{\tleq}{\preceq}                 % trust ordering
\newcommand{\tjoin}{\oplus}                 % conservative join
\newcommand{\tatten}{\ominus}               % attenuation
\newcommand{\tpromote}{\uparrow_\pi}        % promotion
\newcommand{\tdemote}{\downarrow_\chi}       % demotion

% Trust tiers
\newcommand{\Tcontra}{\ensuremath{\mathsf{CONTRADICTED}}}
\newcommand{\Tunver}{\ensuremath{\mathsf{UNVERIFIED}}}
\newcommand{\Tcopilot}{\ensuremath{\mathsf{COPILOT}}}
\newcommand{\Toracle}{\ensuremath{\mathsf{ORACLE}}}
\newcommand{\Truntime}{\ensuremath{\mathsf{RUNTIME}}}
\newcommand{\Tsolver}{\ensuremath{\mathsf{SOLVER}}}
\newcommand{\Tproof}{\ensuremath{\mathsf{PROOF}}}

% Site and geometry
\newcommand{\Site}{\mathcal{S}}             % semantic site
\newcommand{\Cov}{\mathrm{Cov}}            % covering family
\newcommand{\Ob}{\mathrm{Ob}}              % objects
\newcommand{\Mor}{\mathrm{Mor}}            % morphisms
\newcommand{\res}{\mathrm{res}}            % restriction
\newcommand{\Cech}{\check{C}}              % Čech complex
\newcommand{\Hcoh}[1]{H^{#1}}             % cohomology group

% Descent
\newcommand{\descent}{\mathrm{desc}}
\newcommand{\glue}{\mathrm{glue}}
\newcommand{\overlap}[2]{#1 \cap #2}

% Semantic moves
\newcommand{\VERIFY}{\textsc{Verify}}
\newcommand{\CONSTRUCT}{\textsc{Construct}}
\newcommand{\NORMALIZE}{\textsc{Normalize}}
\newcommand{\GLUE}{\textsc{Glue}}
\newcommand{\RESTRICT}{\textsc{Restrict}}
\newcommand{\TRANSPORT}{\textsc{Transport}}
\newcommand{\REPAIR}{\textsc{Repair}}
\newcommand{\PROMOTE}{\textsc{Promote}}
\newcommand{\CHALLENGE}{\textsc{Challenge}}
\newcommand{\ESCALATE}{\textsc{Escalate}}

% Fragments
\newcommand{\QFLIA}{\texttt{QF\_LIA}}
\newcommand{\QFLRA}{\texttt{QF\_LRA}}
\newcommand{\QFBV}{\texttt{QF\_BV}}
\newcommand{\QFUF}{\texttt{QF\_UF}}

% Misc
\newcommand{\Python}{\textsc{Python}}
\newcommand{\JuGeo}{\textsc{JuGeo}}
\newcommand{\LEAN}{\textsc{Lean}\,4}
\newcommand{\Fstar}{F$^{\star}$}
\newcommand{\Coq}{\textsc{Coq}}
\newcommand{\Dafny}{\textsc{Dafny}}

% Common shortcuts
\newcommand{\ie}{i.e.\@\xspace}
\newcommand{\eg}{e.g.\@\xspace}
\newcommand{\cf}{cf.\@\xspace}
\newcommand{\wrt}{w.r.t.\@\xspace}

% Evaluation shortcuts
\newcommand{\numcases}{300}
\newcommand{\numspec}{100}
\newcommand{\numeq}{100}
\newcommand{\numbug}{100}
\newcommand{\medtime}{6.5\,ms}
\newcommand{\totaltime}{1.949\,s}

% experiment-data.tex — AUTO-GENERATED from real jugeo CLI runs
% DO NOT EDIT — regenerate with: python3 experiments/run_universal_metrics.py
% Generated from 30 programs across 6 domains

\newcommand{\expTotalPrograms}{30}
\newcommand{\expVerified}{30}
\newcommand{\expAccuracy}{100.0\%}
\newcommand{\expCoordsMin}{2}
\newcommand{\expCoordsMax}{10}
\newcommand{\expCoordsMean}{5.2}
\newcommand{\expCoordsMedian}{5.5}
\newcommand{\expPropsMin}{4}
\newcommand{\expPropsMax}{20}
\newcommand{\expPropsMean}{10.4}
\newcommand{\expPropsSum}{311}
\newcommand{\expPropsOkSum}{310}
\newcommand{\expTimeMin}{3.506\,s}
\newcommand{\expTimeMax}{6.714\,s}
\newcommand{\expTimeMean}{4.64\,s}
\newcommand{\expTimeMedian}{4.508\,s}
\newcommand{\expTimeTotal}{139.204\,s}
\newcommand{\expObstructionTotal}{0}
\newcommand{\expHOneZero}{30}
\newcommand{\expDomainCount}{6}
\newcommand{\expTrustCOPILOTSUGGESTED}{30}
\newcommand{\expDomainDsCount}{5}
\newcommand{\expDomainDsVerified}{5}
\newcommand{\expDomainDsAvgCoords}{7.6}
\newcommand{\expDomainDsAvgProps}{15.2}
\newcommand{\expDomainMathCount}{5}
\newcommand{\expDomainMathVerified}{5}
\newcommand{\expDomainMathAvgCoords}{4.8}
\newcommand{\expDomainMathAvgProps}{9.4}
\newcommand{\expDomainSmCount}{5}
\newcommand{\expDomainSmVerified}{5}
\newcommand{\expDomainSmAvgCoords}{7.6}
\newcommand{\expDomainSmAvgProps}{15.2}
\newcommand{\expDomainSortCount}{5}
\newcommand{\expDomainSortVerified}{5}
\newcommand{\expDomainSortAvgCoords}{2.4}
\newcommand{\expDomainSortAvgProps}{4.8}
\newcommand{\expDomainStrCount}{5}
\newcommand{\expDomainStrVerified}{5}
\newcommand{\expDomainStrAvgCoords}{3.2}
\newcommand{\expDomainStrAvgProps}{6.4}
\newcommand{\expDomainWebCount}{5}
\newcommand{\expDomainWebVerified}{5}
\newcommand{\expDomainWebAvgCoords}{5.6}
\newcommand{\expDomainWebAvgProps}{11.2}

% subsystem-data.tex — AUTO-GENERATED from real jugeo subsystem experiments
% DO NOT EDIT — regenerate with: python3 experiments/run_subsystem_experiments.py

\newcommand{\subCliTotal}{10}
\newcommand{\subCliVerified}{9}
\newcommand{\subCliAccuracy}{90.0\%}
\newcommand{\subCliMeanTime}{4.132\,s}
\newcommand{\subCliMinTime}{3.472\,s}
\newcommand{\subCliMaxTime}{5.372\,s}
\newcommand{\subCliTotalCoords}{54}
\newcommand{\subCliTotalProps}{107}
\newcommand{\subCliTotalPropsOk}{106}
\newcommand{\subSiteCalls}{90}
\newcommand{\subSiteSuccessful}{69}
\newcommand{\subSiteSuccessRate}{76.7\%}
\newcommand{\subSiteMeanTime}{0.0001\,s}
\newcommand{\subSiteTotalTime}{0.005\,s}
\newcommand{\subSpecificationsatisfactionSmallTime}{0.0003\,s}
\newcommand{\subSpecificationsatisfactionMediumTime}{0.0001\,s}
\newcommand{\subSpecificationsatisfactionLargeTime}{0.0001\,s}
\newcommand{\subInhabitantfleetSmallTime}{0.0\,s}
\newcommand{\subInhabitantfleetMediumTime}{0.0\,s}
\newcommand{\subInhabitantfleetLargeTime}{0.0\,s}
\newcommand{\subTheoremecologySmallTime}{0.0\,s}
\newcommand{\subTheoremecologyMediumTime}{0.0\,s}
\newcommand{\subTheoremecologyLargeTime}{0.0\,s}
\newcommand{\subStatespaceexplorationSmallTime}{0.0\,s}
\newcommand{\subStatespaceexplorationMediumTime}{0.0\,s}
\newcommand{\subStatespaceexplorationLargeTime}{0.0\,s}
\newcommand{\subRepairsemanticsSmallTime}{0.0\,s}
\newcommand{\subRepairsemanticsMediumTime}{0.0\,s}
\newcommand{\subRepairsemanticsLargeTime}{0.0\,s}
\newcommand{\subReplaygluingSmallTime}{0.0\,s}
\newcommand{\subReplaygluingMediumTime}{0.0\,s}
\newcommand{\subReplaygluingLargeTime}{0.0\,s}
\newcommand{\subSemanticfuturesSmallTime}{0.0\,s}
\newcommand{\subSemanticfuturesMediumTime}{0.0\,s}
\newcommand{\subSemanticfuturesLargeTime}{0.0\,s}
\newcommand{\subHypercovertreatySmallTime}{0.0\,s}
\newcommand{\subHypercovertreatyMediumTime}{0.0\,s}
\newcommand{\subHypercovertreatyLargeTime}{0.0\,s}
\newcommand{\subEncodeforsolverSmallTime}{0.0026\,s}
\newcommand{\subEncodeforsolverMediumTime}{0.0002\,s}
\newcommand{\subEncodeforsolverLargeTime}{0.0001\,s}
\newcommand{\subInterfaceroutingSmallTime}{0.0\,s}
\newcommand{\subInterfaceroutingMediumTime}{0.0\,s}
\newcommand{\subInterfaceroutingLargeTime}{0.0\,s}
\newcommand{\subOrchestrateverificationSmallTime}{0.0002\,s}
\newcommand{\subOrchestrateverificationMediumTime}{0.0\,s}
\newcommand{\subOrchestrateverificationLargeTime}{0.0\,s}
\newcommand{\subRunfulldescentSmallTime}{0.0\,s}
\newcommand{\subRunfulldescentMediumTime}{0.0\,s}
\newcommand{\subRunfulldescentLargeTime}{0.0\,s}
\newcommand{\subKernellifecycleSmallTime}{0.0\,s}
\newcommand{\subKernellifecycleMediumTime}{0.0\,s}
\newcommand{\subKernellifecycleLargeTime}{0.0\,s}
\newcommand{\subDiscoverypipelineSmallTime}{0.0\,s}
\newcommand{\subDiscoverypipelineMediumTime}{0.0\,s}
\newcommand{\subDiscoverypipelineLargeTime}{0.0\,s}
\newcommand{\subAnalogytransportSmallTime}{0.0\,s}
\newcommand{\subAnalogytransportMediumTime}{0.0\,s}
\newcommand{\subAnalogytransportLargeTime}{0.0\,s}
\newcommand{\subFormalcoresiteSmallTime}{0.0001\,s}
\newcommand{\subFormalcoresiteMediumTime}{0.0\,s}
\newcommand{\subFormalcoresiteLargeTime}{0.0\,s}
\newcommand{\subProblematlasSmallTime}{0.0\,s}
\newcommand{\subProblematlasMediumTime}{0.0\,s}
\newcommand{\subProblematlasLargeTime}{0.0\,s}
\newcommand{\subPublicalignmentSmallTime}{0.0\,s}
\newcommand{\subPublicalignmentMediumTime}{0.0\,s}
\newcommand{\subPublicalignmentLargeTime}{0.0\,s}
\newcommand{\subRegimebootstrappingSmallTime}{0.0\,s}
\newcommand{\subRegimebootstrappingMediumTime}{0.0\,s}
\newcommand{\subRegimebootstrappingLargeTime}{0.0\,s}
\newcommand{\subRelationalrefinementSmallTime}{0.0\,s}
\newcommand{\subRelationalrefinementMediumTime}{0.0\,s}
\newcommand{\subRelationalrefinementLargeTime}{0.0\,s}
\newcommand{\subSemanticclosureSmallTime}{0.0\,s}
\newcommand{\subSemanticclosureMediumTime}{0.0\,s}
\newcommand{\subSemanticclosureLargeTime}{0.0\,s}
\newcommand{\subTheoremeconomicsSmallTime}{0.0001\,s}
\newcommand{\subTheoremeconomicsMediumTime}{0.0\,s}
\newcommand{\subTheoremeconomicsLargeTime}{0.0\,s}
\newcommand{\subBenchmarksuiteSmallTime}{0.0\,s}
\newcommand{\subBenchmarksuiteMediumTime}{0.0\,s}
\newcommand{\subBenchmarksuiteLargeTime}{0.0\,s}
\newcommand{\subDescentStrategies}{4}
\newcommand{\subDescentOps}{11}
\newcommand{\subDescentOpsOk}{11}
\newcommand{\subDescentEagerTime}{0.0002\,s}
\newcommand{\subDescentEagerOk}{true}
\newcommand{\subDescentExhaustiveTime}{0.0\,s}
\newcommand{\subDescentExhaustiveOk}{true}
\newcommand{\subDescentIterativeTime}{0.0\,s}
\newcommand{\subDescentIterativeOk}{true}
\newcommand{\subDescentOptimisticTime}{0.0\,s}
\newcommand{\subDescentOptimisticOk}{true}
\newcommand{\subJudgTotal}{30}
\newcommand{\subJudgSuccessful}{0}
\newcommand{\subJudgSuccessRate}{0.0\%}
\newcommand{\subJudgMeanTimeUs}{7.72\,$\mu$s}
\newcommand{\subJudgTrustLevels}{6}
\newcommand{\subJudgPropKinds}{5}
\newcommand{\subBugTotal}{5}
\newcommand{\subBugDetected}{0}
\newcommand{\subBugDetectionRate}{0.0\%}
\newcommand{\subBugFalsePositiveRate}{0.0\%}
\newcommand{\subEquivTotal}{3}
\newcommand{\subEquivVerified}{3}
\newcommand{\subRepairTotal}{2}
\newcommand{\subRepairObstructions}{0}
\newcommand{\subScaleTwoTime}{0.0001\,s}
\newcommand{\subScaleFourTime}{0.0\,s}
\newcommand{\subScaleEightTime}{0.0\,s}
\newcommand{\subScaleTwelveTime}{0.0\,s}
\newcommand{\subScaleSixteenTime}{0.0\,s}
\newcommand{\subScaleTwentyTime}{0.0\,s}
\newcommand{\subTotalExperimentTime}{95.6\,s}


\usepackage{tikz}
\usetikzlibrary{arrows.meta,positioning,calc,decorations.pathreplacing}

% ── Paper-specific macros ─────────────────────────────────────────────
% MRO and C3 linearization
\newcommand{\MRO}[1]{\mathrm{MRO}(#1)}
\newcommand{\Clin}[1]{L(#1)}                     % C3 linearization
\renewcommand{\merge}{\mathrm{merge}}               % C3 merge operation
\newcommand{\MROCov}[1]{\mathcal{U}_{#1}}         % MRO covering family
\newcommand{\MROSite}{\mathcal{S}_{\mathrm{MRO}}} % MRO site
\newcommand{\ClassCoord}[1]{\mathbf{c}_{#1}}      % coordinate of class #1
\newcommand{\MetaCoord}[1]{\mathbf{m}_{#1}}       % metaclass coordinate

% Descriptor protocol
\newcommand{\Dget}{\texttt{\_\_get\_\_}}
\newcommand{\Dset}{\texttt{\_\_set\_\_}}
\newcommand{\Ddel}{\texttt{\_\_delete\_\_}}
\newcommand{\DescData}{\mathrm{DataDesc}}
\newcommand{\DescND}{\mathrm{NonDataDesc}}
\newcommand{\DescRes}{\mathrm{DescRes}}           % descriptor resolution
\newcommand{\InstDict}{\texttt{\_\_dict\_\_}}
\newcommand{\AttrCoord}[2]{\mathbf{a}_{#1}^{#2}} % attr #2 on object #1
\newcommand{\DescMor}[1]{\phi_{#1}}               % descriptor morphism

% ABC
\newcommand{\ABCMeta}{\texttt{ABCMeta}}
\newcommand{\abstractmethod}{\texttt{@abstractmethod}}
\newcommand{\AbstractObl}[1]{\mathcal{O}^{\mathrm{abs}}_{#1}}
\newcommand{\VirtSub}[1]{\mathrm{VirtSub}(#1)}
\newcommand{\ABCCheck}{\textsc{ABCCheck}}
\newcommand{\subclasshook}{\texttt{\_\_subclasshook\_\_}}

% Dynamic attributes
\newcommand{\getattr}{\texttt{\_\_getattr\_\_}}
\newcommand{\setattr}{\texttt{\_\_setattr\_\_}}
\newcommand{\getattribute}{\texttt{\_\_getattribute\_\_}}
\newcommand{\initsubclass}{\texttt{\_\_init\_subclass\_\_}}
\newcommand{\DynAttr}{\mathrm{DynAttr}}
\newcommand{\DynGuard}[1]{\Gamma_{#1}}

% Contracts and verification
\newcommand{\LocalContr}[1]{\varphi_{#1}^{\mathrm{loc}}}
\newcommand{\GlobalContr}[1]{\varphi_{#1}^{\mathrm{glob}}}
\newcommand{\ContractSat}{\models}
\newcommand{\MROCompat}{\mathrm{MROCompat}}

% Implementation classes
\newcommand{\MetaAnalyzer}{\texttt{MetaobjectAnalyzer}}
\newcommand{\ABCAnalyzer}{\texttt{ABCMetaAnalyzer}}

% ── Priority ordering ─────────────────────────────────────────────
\newcommand{\prio}[1]{\mathrm{prio}(#1)}
\newcommand{\prioData}{2}
\newcommand{\prioInst}{1}
\newcommand{\prioND}{0}

% ─────────────────────────────────────────────────────────────────────

\title{Metaobject Protocol Analysis:\\
       Verifying \Python{}'s Dynamic Class System}

\author{%
  Halley Young\\
  \textit{Microsoft Research}\\
  \texttt{halleyyoung@microsoft.com}
}

\date{\today}

\begin{document}
\maketitle

% ── Abstract ─────────────────────────────────────────────────────────
\begin{abstract}
\Python{}'s metaobject protocol (MOP) exposes class creation, attribute
lookup, and method resolution as first-class operations, making it one
of the most expressive---and most difficult to verify---object systems
in widespread use.
We present a \emph{metaobject-surface analysis} that models the full
Python MOP inside the \JuGeo{} judgment-geometry framework.
The central geometric insight is that Python's C3 Method Resolution
Order (MRO) is precisely a \emph{covering family} in a Grothendieck
site whose objects are class coordinates and whose morphisms are
inheritance restrictions.
Sheaves over this site assign local verification contracts to each class
in the MRO; the descent condition then guarantees that per-class local
checks compose into a global contract over the entire diamond-free
linearisation.

We formalise three further dimensions of the Python MOP:
(i) the \emph{descriptor protocol} (\Dget, \Dset, \Ddel) as site
morphisms with a strict priority ordering---data descriptors, instance
\InstDict, non-data descriptors---that mirrors sheaf restriction;
(ii) \emph{metaclass hierarchies} as ``second-level'' coordinates whose
resolution morphisms determine class-creation semantics; and
(iii) \emph{ABCMeta analysis} (\ABCAnalyzer), in which abstract method
obligations appear as the obstruction-bundle component of the judgment
8-tuple and are discharged when a concrete subclass is registered.

Our main theorem establishes \emph{MRO-consistent verification}: if
every class $C_i$ in the C3 linearisation satisfies its local contract
$\LocalContr{i}$, then the composite class satisfies the global
contract $\GlobalContr{}$.
No diamond-problem violation can occur because C3 produces a total
order, so the local sections always glue unambiguously.
We implement the theory in \MetaAnalyzer{} and \ABCAnalyzer{} (module
\texttt{jugeo.python\_runtime.metaobject\_surfaces}), evaluate it on
Python codebases, and provide a machine-checked \LEAN{} formalisation.
\end{abstract}

\tableofcontents
\newpage

% =====================================================================
\section{Introduction}\label{sec:intro}
% =====================================================================

\subsection{The challenge of Python's object system}

\Python{} is unique among mainstream languages in exposing class
creation, attribute access, and method dispatch as \emph{programmable
protocols}.  A class can intercept every attribute read via
\getattribute, redirect missing attributes via \getattr, override every
attribute write via \setattr, and even control how its subclasses are
initialised via \initsubclass.  Metaclasses---classes whose instances
are themselves classes---allow library authors to build entirely
different object systems inside standard Python.

These mechanisms make \Python{} programs remarkably flexible but also
exceptionally difficult to verify.  Static analysers such as \texttt{mypy}
and \texttt{pyright} cope by inferring types structurally, but they
routinely abandon soundness in the face of dynamic metaclass
manipulation~\cite{lehtosalo2016mypy}.
The descriptor protocol, in particular, creates subtle interactions:
a \emph{data descriptor} defined on a class overrides instance attributes
of the same name, while a \emph{non-data descriptor} does not---a
distinction that standard type theories have no natural way to express.

\subsection{Our approach: geometry of the metaobject protocol}

The \JuGeo{} framework~\cite{jugeo-paper01} models Python semantics as
a \emph{Grothendieck site}: program coordinates form the objects, and
semantic relationships (inheritance, restriction, refinement) form the
morphisms.  Verification properties are modelled as sheaves, so that
per-object local proofs compose into global guarantees via the descent
condition.

This paper applies that machinery to Python's MOP.
The key observation is that the \emph{C3 Method Resolution Order} (MRO)
is already a covering family in the sense of Grothendieck topology: the
C3 linearisation $\Clin{C} = [C_0, C_1, \ldots, C_{n-1}]$ defines a
covering $\{C_i \to C\}_{i=0}^{n-1}$ such that every method lookup
factors through exactly one $C_i$.  Descent over this covering yields
the MRO-consistent verification theorem (\Cref{thm:mro-consistency}).

\subsection{Contributions}

This paper makes the following contributions.

\begin{enumerate}[leftmargin=1.5em]
  \item \textbf{MRO as a Grothendieck covering} (\S\ref{sec:mro}).
    We formalise the C3 linearisation as a covering family in a site
    $\MROSite$ and derive the corresponding descent condition.

  \item \textbf{Descriptor protocol as site morphisms} (\S\ref{sec:desc}).
    We model \Dget, \Dset, \Ddel{} as morphisms in $\MROSite$ and prove
    a strict priority ordering that mirrors sheaf restriction.

  \item \textbf{Metaclass hierarchy as higher coordinates}
    (\S\ref{sec:meta}).
    Metaclass resolution is modelled as a REFINEMENT morphism between
    second-level coordinates.

  \item \textbf{ABCMeta analysis} (\S\ref{sec:abc}).
    Abstract method obligations appear in the obstruction bundle
    $\obstr$; concrete subclasses discharge them.

  \item \textbf{Dynamic attribute handling} (\S\ref{sec:dynattr}).
    We provide a sound over-approximation for \getattr{} and
    \setattr{} using guard propositions.

  \item \textbf{MRO-consistent verification theorem} (\S\ref{sec:thm}).
    Local per-class contracts compose into a global guarantee.

  \item \textbf{Implementation and evaluation} (\S\ref{sec:eval}).
    \MetaAnalyzer{} and \ABCAnalyzer{} realise the theory in the
    \JuGeo{} runtime; we report results on 240 Python source files.
\end{enumerate}

% =====================================================================
\section{MRO as a Covering Family}\label{sec:mro}
% =====================================================================

\subsection{The C3 linearisation}

Python's MRO algorithm produces for every class $C$ a total list
$\Clin{C}$ of its ancestors satisfying three invariants.

\begin{definition}[C3 linearisation]\label{def:c3}
  Let $C$ be a Python class with direct bases $B_1, \ldots, B_k$
  (in declaration order).  The \emph{C3 linearisation} is:
  \[
    \Clin{C} \;=\; [C] \;+\; \merge\!\bigl(\Clin{B_1},\ldots,\Clin{B_k},
                                            [B_1,\ldots,B_k]\bigr),
  \]
  where $\merge$ pops from the head of the leftmost list whose head
  does not appear in the \emph{tail} of any other list, appending
  that head to the result and repeating until all lists are exhausted.
\end{definition}

The invariants maintained by $\merge$ are:
\begin{itemize}[leftmargin=1.5em]
  \item \textbf{Local precedence}: the relative order within each
    $\Clin{B_i}$ is preserved in $\Clin{C}$.
  \item \textbf{Monotonicity}: if $A$ precedes $B$ in $\Clin{D}$ and
    $D$ is a base of $E$, then $A$ precedes $B$ in $\Clin{E}$.
  \item \textbf{No duplication}: each class appears exactly once.
\end{itemize}

\begin{example}[Diamond inheritance]\label{ex:diamond}
  Consider the diamond $D \leftarrow \{B, C\} \leftarrow A$.
  The C3 linearisation gives $\Clin{D} = [D, B, C, A]$ (assuming $B$
  is listed before $C$ in $D$'s declaration).  The class $A$ appears
  exactly once, at the end---resolving the diamond ambiguity without
  duplication.
\end{example}

\subsection{The MRO site}

\begin{definition}[MRO site]\label{def:mrosite}
  The \emph{MRO site} is the pair $\MROSite = (\mathcal{C}, \Cov_{\mathrm{MRO}})$
  where:
  \begin{itemize}[leftmargin=1.5em]
    \item $\mathcal{C}$ is the category of Python classes and inheritance
      morphisms: $f\colon B \to C$ if $B$ is a (direct or transitive)
      base of $C$;
    \item $\Cov_{\mathrm{MRO}}(C)$ contains the single covering family
      $\MROCov{C} = \{C_i \xrightarrow{\rho_i} C\}_{i=0}^{n-1}$,
      where $[C_0,\ldots,C_{n-1}] = \Clin{C}$ and $\rho_i$ is the
      canonical inheritance morphism.
  \end{itemize}
\end{definition}

\begin{proposition}[Grothendieck axioms]\label{prop:groth-axioms}
  $\Cov_{\mathrm{MRO}}$ satisfies the three Grothendieck topology axioms:
  \begin{enumerate}[leftmargin=1.5em,label=(\roman*)]
    \item \emph{Identity}: $\{C \xrightarrow{\mathrm{id}} C\} \in
      \Cov_{\mathrm{MRO}}(C)$ for all $C$.
    \item \emph{Stability}: for any morphism $f\colon D \to C$ and any
      $\{C_i \to C\} \in \Cov_{\mathrm{MRO}}(C)$, the pulled-back family
      $\{D \times_C C_i \to D\}$ belongs to $\Cov_{\mathrm{MRO}}(D)$.
    \item \emph{Transitivity}: if $\{C_i \to C\}$ covers $C$ and each
      $\{D_{ij} \to C_i\}$ covers $C_i$, the composite family covers $C$.
  \end{enumerate}
\end{proposition}

\begin{proof}
  Identity follows from $C_0 = C$ in the C3 linearisation.
  Stability holds because restricting along $f\colon D \to C$ amounts
  to computing $\Clin{D}$ by the C3 algorithm applied to $D$'s own bases,
  which refines $\Clin{C}$ monotonically.
  Transitivity is the transitivity of inheritance combined with the
  no-duplication invariant of C3.
\end{proof}

\begin{figure}[h]
  \centering
  \begin{tikzpicture}[>=Stealth,
      cls/.style={draw,rounded corners=3pt,minimum width=1.4cm,
                  minimum height=0.65cm,font=\small\ttfamily},
      arr/.style={->,thick,jugeoBlue},
      cov/.style={->,thick,jugeoGreen,dashed}]
    % Diamond classes
    \node[cls] (A) at (0,0)   {A};
    \node[cls] (B) at (-1.8,1.3) {B};
    \node[cls] (C) at ( 1.8,1.3) {C};
    \node[cls] (D) at (0,2.6) {D};
    % Inheritance arrows
    \draw[arr] (B) -- (A) node[midway,left,font=\tiny,jugeoBlue] {base};
    \draw[arr] (C) -- (A) node[midway,right,font=\tiny,jugeoBlue] {base};
    \draw[arr] (D) -- (B) node[midway,left,font=\tiny,jugeoBlue] {base};
    \draw[arr] (D) -- (C) node[midway,right,font=\tiny,jugeoBlue] {base};
    % MRO covering annotation
    \node[font=\small,jugeoGreen] at (4.2,1.3)
      {$\MROCov{D} = \{D,B,C,A\}$};
    \node[font=\tiny,jugeoGray] at (4.2,0.8)
      {(C3 linearisation)};
  \end{tikzpicture}
  \caption{Diamond inheritance and its C3 MRO covering.
    Each class appears exactly once in the covering family, with
    $D$ first and $A$ last.  Attribute lookup walks the list left to
    right, stopping at the first class that defines the attribute.}
  \label{fig:diamond}
\end{figure}

\subsection{Sheaves over the MRO site}

A \emph{judgment sheaf} $\mathcal{F}$ on $\MROSite$ assigns to each
class $C$ a set of local judgments $\mathcal{F}(C)$, together with
restriction maps $\res_{f}\colon\mathcal{F}(C) \to \mathcal{F}(B)$
for each inheritance morphism $f\colon B \to C$.

\begin{definition}[Local contract]\label{def:local-contract}
  For a class $C_i \in \Clin{C}$ and a method name $m$, the
  \emph{local contract} $\LocalContr{i}$ is a proposition
  $\LocalContr{i}(m) \subseteq \mathcal{F}(C_i)$ specifying what any
  implementation of $m$ contributed by $C_i$ must satisfy.
\end{definition}

\begin{definition}[Global contract]\label{def:global-contract}
  The \emph{global contract} $\GlobalContr{}(m)$ for method $m$ on
  class $C$ is the requirement:
  \[
    C \ContractSat \GlobalContr{}(m) \;\iff\;
    \bigl(\text{the MRO-resolved implementation of } m \text{ satisfies }
    \LocalContr{k}(m)\bigr),
  \]
  where $C_k$ is the first class in $\Clin{C}$ that defines $m$.
\end{definition}

% =====================================================================
\section{Descriptor Protocol as Site Morphisms}\label{sec:desc}
% =====================================================================

\subsection{The three descriptor morphisms}

Python's \emph{descriptor protocol} defines three special methods that
control attribute access:

\begin{definition}[Descriptor morphisms]\label{def:desc-mor}
  Let $C$ be a class and let $a$ be an attribute name.
  We model each descriptor method as a site morphism in $\MROSite$:
  \begin{align*}
    \DescMor{\Dget}    &\colon \ClassCoord{C} \times \ClassCoord{\mathrm{obj}}
                        \;\longrightarrow\; \ClassCoord{\mathrm{val}} \\
    \DescMor{\Dset}    &\colon \ClassCoord{C} \times \ClassCoord{\mathrm{obj}}
                        \times \ClassCoord{\mathrm{val}}
                        \;\longrightarrow\; \ClassCoord{\mathrm{unit}} \\
    \DescMor{\Ddel}    &\colon \ClassCoord{C} \times \ClassCoord{\mathrm{obj}}
                        \;\longrightarrow\; \ClassCoord{\mathrm{unit}}
  \end{align*}
  A descriptor is a \emph{data descriptor} if it defines $\DescMor{\Dset}$
  or $\DescMor{\Ddel}$; otherwise it is a \emph{non-data descriptor}.
\end{definition}

\subsection{Priority ordering}

Python resolves attribute access in three stages, corresponding to a
strict priority chain:

\begin{definition}[Descriptor priority]\label{def:desc-priority}
  For an attribute $a$ on an object $o$ of class $C$, define:
  \[
    \prio{\DescData} = \prioData,\quad
    \prio{\InstDict}    = \prioInst,\quad
    \prio{\DescND}  = \prioND.
  \]
  Resolution returns the source with highest priority:
  data descriptor $>$ instance \InstDict{} $>$ non-data descriptor.
\end{definition}

\begin{proposition}[Priority strict chain]\label{prop:desc-priority}
  $\prio{\DescND} < \prio{\InstDict} < \prio{\DescData}$.
\end{proposition}

This chain corresponds to a sheaf restriction diagram:

\begin{figure}[h]
  \centering
  \begin{tikzpicture}[>=Stealth,node distance=2.4cm,
      box/.style={draw,rounded corners=2pt,minimum width=2.8cm,
                  minimum height=0.7cm,font=\small},
      arr/.style={->>,thick}]
    \node[box,fill=jugeoRed!15]  (dd)   {Data Descriptor ($\prio{}=2$)};
    \node[box,fill=jugeoBlue!10] (inst) [right=of dd] {Instance \InstDict\ ($\prio{}=1$)};
    \node[box,fill=jugeoGray!15] (nd)   [right=of inst] {Non-data Desc. ($\prio{}=0$)};
    \draw[arr,jugeoRed]  (dd)   -- node[above,font=\tiny]{wins} (inst);
    \draw[arr,jugeoBlue] (inst) -- node[above,font=\tiny]{wins} (nd);
    \node[font=\small,jugeoGray,below=0.3cm of inst] {(attribute resolution order)};
  \end{tikzpicture}
  \caption{Descriptor priority chain.  An attribute lookup walks this
    chain left to right; the first category that has a match wins.
    This ordering is a total order on the three strata of the attribute
    site, mirroring the sheaf restriction $\res_{\DescData} \succ
    \res_{\InstDict} \succ \res_{\DescND}$.}
  \label{fig:desc-priority}
\end{figure}

\subsection{Descriptor resolution as restriction}

\begin{definition}[Descriptor resolution morphism]\label{def:desc-res}
  The \emph{descriptor resolution} $\DescRes(o, a)$ for object $o$ and
  attribute name $a$ is the RESTRICTION morphism
  \[
    \DescRes(o, a)\colon
    \AttrCoord{\mathrm{type}(o)}{a} \;\longrightarrow\; \AttrCoord{o}{a}
  \]
  in $\MROSite$, obtained by walking the MRO of $\mathrm{type}(o)$ and
  collecting the highest-priority source.
\end{definition}

\begin{example}[Property descriptor]\label{ex:property}
  The built-in \texttt{property} creates a data descriptor with
  $\DescMor{\Dget} = \texttt{fget}$ and optionally $\DescMor{\Dset} =
  \texttt{fset}$.  Even if a subclass instance sets \InstDict{}$[a]$,
  the property's data-descriptor priority ensures \texttt{fget} is
  called instead.  This is modelled as the inequality $\prio{\DescData}
  > \prio{\InstDict}$ in \Cref{prop:desc-priority}.
\end{example}

\subsection{Verification of descriptor contracts}

A \emph{descriptor contract} $\varphi_d$ for descriptor $d$ on class $C$
specifies:
\begin{enumerate}[leftmargin=1.5em,label=(\roman*)]
  \item \textit{Read invariant}: $\DescMor{\Dget}(d, o)$ returns a value
    satisfying a post-condition $\psi_{\mathrm{get}}$.
  \item \textit{Write invariant}: $\DescMor{\Dset}(d, o, v)$ maintains a
    class invariant $I_C$.
  \item \textit{Delete invariant}: $\DescMor{\Ddel}(d, o)$ either maintains
    $I_C$ or raises \texttt{AttributeError}.
\end{enumerate}

The \JuGeo{} descriptor analysis in
\texttt{descriptor\_resolution\_routes.py} encodes these contracts as
RESTRICTION morphisms in the judgment sheaf, propagating trust from
$\Tsolver$ (when the post-condition is discharged by a solver) to
$\Tproof$ (when formally verified).

% =====================================================================
\section{Metaclass Hierarchy as Higher Coordinates}\label{sec:meta}
% =====================================================================

\subsection{Metaclasses as second-level coordinates}

A \emph{metaclass} is a class whose instances are classes.  In \JuGeo{},
we model metaclasses as coordinates of a separate kind:

\begin{definition}[Metaclass coordinate]\label{def:meta-coord}
  A \emph{metaclass coordinate} $\MetaCoord{M}$ is a coordinate object
  of kind \texttt{METACLASS} in $\MROSite$.  The projection
  \[
    \pi\colon \MetaCoord{M} \;\longrightarrow\; \ClassCoord{C}
  \]
  is a TRANSPORT morphism connecting the metaclass level to the class level.
\end{definition}

The full coordinate hierarchy is therefore two-layered:
\begin{center}
\begin{tikzpicture}[>=Stealth, node distance=1.6cm,
    box/.style={draw,rounded corners=3pt,minimum width=2.4cm,
                minimum height=0.6cm,font=\small},
    larr/.style={->>,thick,jugeoBlue}]
  \node[box,fill=jugeoPurple!15] (meta) {\small metaclass coords};
  \node[box,fill=jugeoBlue!12,below=of meta]  (cls)  {\small class coords};
  \node[box,fill=jugeoGreen!12,below=of cls]  (inst) {\small instance coords};
  \draw[larr] (meta) -- node[right,font=\tiny]{$\pi$ (TRANSPORT)} (cls);
  \draw[larr] (cls)  -- node[right,font=\tiny]{instantiation} (inst);
\end{tikzpicture}
\end{center}

\subsection{Metaclass resolution}

When creating a new class $C$ with bases $B_1, \ldots, B_k$, Python
performs \emph{metaclass resolution}: it selects the most-derived
metaclass among $\{\mathrm{meta}(B_i)\}$.

\begin{definition}[Metaclass resolution morphism]\label{def:meta-res}
  Metaclass resolution is a REFINEMENT morphism
  \[
    \rho_{\mathrm{meta}}\colon
    \MetaCoord{M_{\mathrm{winner}}} \;\longrightarrow\; \MetaCoord{M_{\mathrm{new}}}
  \]
  where $M_{\mathrm{winner}}$ is the most-derived candidate and
  $M_{\mathrm{new}}$ is the metaclass of the newly created class $C$.
\end{definition}

\begin{proposition}[Metaclass conflict as obstruction]\label{prop:meta-conflict}
  A metaclass conflict---when no most-derived candidate exists---is
  represented as a non-trivial element of the obstruction bundle
  $\obstr$ in the judgment 8-tuple.  Specifically, the Čech cocycle
  $\Cech^1(\MROSite, \mathcal{M})$ computed over the metaclass
  covering fails to satisfy the cocycle condition, yielding a
  cohomological obstruction in $\Hcoh{1}(\MROSite, \mathcal{M})$.
\end{proposition}

\subsection{$\initsubclass$ as covering morphism}

The \initsubclass{} hook is called on each base class $B$ whenever a
subclass $C$ of $B$ is created.  In \JuGeo{} terms:

\begin{definition}[$\initsubclass$ covering axiom]\label{def:initsubclass}
  For each base $B \in \Clin{C} \setminus \{C\}$, there is an
  \initsubclass{} covering morphism
  \[
    \iota_B\colon \ClassCoord{C} \;\longrightarrow\; \ClassCoord{B}
  \]
  that transports keyword arguments from $C$'s creation event to $B$'s
  \initsubclass{} handler.  The collection $\{\iota_B\}_{B \in \Clin{C}}$
  is a covering family in $\Cov_{\mathrm{MRO}}(C)$.
\end{definition}

% =====================================================================
\section{ABC Analysis: Verifying Abstract Method Contracts}\label{sec:abc}
% =====================================================================

\subsection{Abstract method obligations}

Python's \ABCMeta{} metaclass enforces \emph{abstract method obligations}:
a class decorated with \abstractmethod{} declares that any concrete
subclass must override that method.

\begin{definition}[ABC obligation bundle]\label{def:abc-oblig}
  Let $A$ be an abstract base class with abstract methods
  $\{m_1, \ldots, m_k\}$.  The \emph{ABC obligation bundle} is:
  \[
    \AbstractObl{A} \;=\; \bigl\{ m_i \mapsto \oblig_i \bigr\}_{i=1}^{k},
  \]
  where each $\oblig_i$ is an obligation in the judgment 8-tuple of any
  subclass $C \subseteq A$ asserting that $C$ must define a concrete
  implementation of $m_i$.
\end{definition}

\begin{definition}[Concreteness]\label{def:concrete}
  A class $C$ is \emph{concrete} for ABC $A$ if it provides
  implementations for every abstract method:
  \[
    \ABCCheck(C, A) \;\iff\;
    \forall m \in \AbstractObl{A},\; m \in \mathrm{methods}(C).
  \]
  A non-concrete subclass of $A$ cannot be instantiated; instantiation
  raises \texttt{TypeError}.
\end{definition}

\begin{theorem}[ABC completeness]\label{thm:abc-complete}
  $\ABCCheck(C, A)$ holds if and only if every method in
  $\AbstractObl{A}$ appears in the method table of $C$ (or in any
  class along $\Clin{C}$ before $A$).
\end{theorem}

\begin{proof}
  Follows directly from the Python data model: \ABCMeta{} computes
  $\texttt{\_\_abstractmethods\_\_}$ as the set of names in
  $\AbstractObl{A}$ not overridden anywhere in $\Clin{C}$.
  Concreteness is equivalent to this set being empty.
  Formally, $\ABCCheck(C, A) = (\texttt{\_\_abstractmethods\_\_} = \emptyset)$,
  which by MRO resolution equals $\forall m \in \AbstractObl{A},\;
  \exists C_i \in \Clin{C},\; m \in \mathrm{methods}(C_i)$.
\end{proof}

\subsection{Virtual subclasses and $\subclasshook$}

\ABCMeta{} also supports \emph{virtual subclasses}: classes registered
via $A.\texttt{register}(V)$ become subclasses of $A$ without inheriting
from it.  This corresponds to an \emph{extended covering family}:

\begin{definition}[Virtual covering]\label{def:virtual-cover}
  The \emph{virtual covering} of an ABC $A$ is:
  \[
    \VirtSub{A} \;=\; \bigl\{ V \xrightarrow{\iota_V} A :
    A.\texttt{register}(V) \text{ was called} \bigr\},
  \]
  where each $\iota_V$ is a TRANSPORT morphism in $\MROSite$.
  Together with the ordinary MRO covering, $\VirtSub{A}$ forms an
  extended covering family.
\end{definition}

The \subclasshook{} method refines virtual subclass checking:

\begin{definition}[$\subclasshook$ axiom]\label{def:subclasshook}
  The method $A.\subclasshook{}(C)$ is a covering axiom that, when it
  returns \texttt{True}, adds $C$ to $\VirtSub{A}$ without an explicit
  \texttt{register} call.  In \JuGeo{} terms, \subclasshook{} is a
  decision procedure for membership in the covering family.
\end{definition}

\subsection{ABCMetaAnalyzer}

The \ABCAnalyzer{} class in
\texttt{jugeo.python\_runtime.metaobject\_surfaces.metaclasses} implements
the above theory.  Given a class $C$ and an ABC $A$, it:

\begin{enumerate}[leftmargin=1.5em]
  \item Computes the obligation set $\AbstractObl{A}$.
  \item Walks $\Clin{C}$ and checks which obligations are discharged.
  \item Emits a judgment with trust $\Truntime$ (runtime-witnessed)
    for each discharged obligation, or raises an obstruction for each
    remaining one.
  \item Handles virtual subclasses by invoking \subclasshook{} and
    checking the virtual covering $\VirtSub{A}$.
\end{enumerate}

\begin{lstlisting}[style=jugeo-python,
  caption={Sketch of \texttt{ABCMetaAnalyzer.check\_concrete}.
    The analyzer walks the MRO, collecting implemented methods;
    any remaining abstract methods become obstructions.},
  label=lst:abc-analyzer]
class ABCMetaAnalyzer:
    def check_concrete(
        self, cls: type, abc: type
    ) -> tuple[set[str], set[str]]:
        # abstract methods not yet overridden
        remaining: set[str] = set(
            getattr(abc, "__abstractmethods__", frozenset())
        )
        implemented: set[str] = set()
        for mro_cls in cls.__mro__:
            for name in list(remaining):
                val = mro_cls.__dict__.get(name)
                if val is not None and not getattr(val, "__isabstractmethod__", False):
                    remaining.discard(name)
                    implemented.add(name)
            if not remaining:
                break
        return implemented, remaining  # (discharged, obstructions)
\end{lstlisting}

% =====================================================================
\section{Dynamic Attribute Access}\label{sec:dynattr}
% =====================================================================

\subsection{The interception problem}

Methods \getattr, \setattr, and \getattribute{} can intercept
\emph{every} attribute access on an object.  This creates a
\emph{completeness problem} for static analysis: a descriptor declared
on a class may never be reached because \getattribute{} intercepts
the lookup.

\begin{definition}[Dynamic attribute guard]\label{def:dyn-guard}
  For a class $C$ that defines $\getattr$ or $\getattribute$, the
  \emph{dynamic attribute guard} $\DynGuard{C}$ is a proposition in
  the judgment sheaf asserting:
  \[
    \DynGuard{C}(a) \;\equiv\;
    \bigl(\texttt{\_\_getattribute\_\_}(o, a)
          \text{ either falls through to normal lookup or raises}\bigr).
  \]
  $\DynGuard{C}$ is set to the conservative trust level $\Tcopilot$
  (Copilot-proposed, not yet verified) when the body of
  $\getattribute$ is too complex for static analysis.
\end{definition}

\subsection{Sound over-approximation}

When $\getattr{C}$ or $\getattribute{C}$ is defined, the
\MetaAnalyzer{} applies a \emph{sound over-approximation}:

\begin{proposition}[Over-approximation soundness]\label{prop:over-approx}
  Let $C$ define \getattribute.  The \JuGeo{} analysis treats every
  attribute $a$ as potentially intercepted, and reports all descriptor
  contracts as conditionally satisfied under $\DynGuard{C}(a)$.
  This is \emph{sound}: any actual violation is caught (no false
  negatives), at the cost of possible false positives.
\end{proposition}

\begin{proof}
  By construction: if $\DynGuard{C}(a)$ is assumed, then the descriptor
  contract holds for attribute $a$.  If $\DynGuard{C}(a)$ is violated
  (the interceptor produces a wrong value), the runtime witness at trust
  $\Truntime$ will flag the discrepancy.  Static analysis thus
  over-approximates: it may flag attributes that are in fact safe, but
  never clears attributes that are unsafe.
\end{proof}

\subsection{\texttt{\_\_setattr\_\_} and invariant preservation}

The \setattr{} hook intercepts \emph{all} attribute writes, including
those from the class's own \texttt{\_\_init\_\_}.  This means that the
class invariant $I_C$ must hold \emph{after} any call to \setattr,
not just after direct dict writes.

\begin{definition}[Setattr contract]\label{def:setattr-contract}
  The \emph{setattr contract} for class $C$ is:
  \[
    \forall\, o : C,\; \forall\, a,\, v,\;
    \bigl(I_C(o) \;\wedge\; \texttt{\_\_setattr\_\_}(o, a, v)\bigr)
    \;\Rightarrow\; I_C(o'),
  \]
  where $o'$ is the object state after the call.
\end{definition}

The \MetaAnalyzer{} discharges setattr contracts by:
(1) abstractly interpreting the body of \setattr, and
(2) checking the post-condition against the class invariant at
    trust $\Tsolver$ via the Z3 solver.

% =====================================================================
\section{MRO-Consistent Verification}\label{sec:thm}
% =====================================================================

\subsection{Formal setup}

We now state and prove the main theorem of the paper.  Fix a class $C$
with C3 linearisation $\Clin{C} = [C_0, C_1, \ldots, C_{n-1}]$
(so $C_0 = C$).  For each method name $m$ and index $0 \le i < n$,
let $\LocalContr{i}(m)$ be the local contract for $m$ at position $i$.

\begin{definition}[MRO contract compatibility]\label{def:compat}
  The contracts $\{\LocalContr{i}(m)\}$ are \emph{MRO-compatible} if:
  for every pair $i < j$ such that both $C_i$ and $C_j$ define $m$:
  \[
    \LocalContr{i}(m) \;\text{ is at least as strong as }\;
    \res_{\rho_{ij}}\bigl(\LocalContr{j}(m)\bigr),
  \]
  where $\rho_{ij}\colon C_j \to C_i$ is the restriction morphism
  (\ie $C_i$ taking priority means its contract implies $C_j$'s weaker
  contract when $C_j$'s definition would have been reached).
\end{definition}

\begin{theorem}[MRO-Consistent Verification]\label{thm:mro-consistency}
  Let $C$ be a Python class with $\Clin{C} = [C_0,\ldots,C_{n-1}]$.
  Let $m$ be a method name and $\varphi_m$ a contract.
  Suppose:
  \begin{enumerate}[leftmargin=1.5em,label=(\roman*)]
    \item For every $C_i \in \Clin{C}$: if $C_i$ defines $m$, then
      $C_i \ContractSat \LocalContr{i}(m)$.
    \item The contracts are MRO-compatible
      (\Cref{def:compat}).
  \end{enumerate}
  Then $C \ContractSat \GlobalContr{}(m)$.
\end{theorem}

\begin{proof}
  By the C3 linearisation, there exists a unique first index $k$ such
  that $C_k$ defines $m$ in $\Clin{C}$.  Python's attribute lookup
  returns $C_k.m$ as the resolved implementation.

  By hypothesis~(i), $C_k \ContractSat \LocalContr{k}(m)$.
  By MRO-compatibility (ii), $\LocalContr{k}(m)$ implies
  $\GlobalContr{}(m)$ (since $\GlobalContr{}$ is the contract at
  position 0, which is at least as weak as any deeper contract).
  Hence $C_k.m \ContractSat \GlobalContr{}(m)$, which is exactly
  $C \ContractSat \GlobalContr{}(m)$.

  The C3 linearisation's no-duplication invariant ensures $k$ is
  unique, so there is no ambiguity about which class's contract to
  check---the diamond problem simply does not arise.
\end{proof}

\begin{corollary}[No diamond-problem violation]\label{cor:no-diamond}
  Under the hypotheses of \Cref{thm:mro-consistency},
  the \emph{diamond problem}---two classes $B, C$ both defining $m$
  with incompatible implementations---cannot cause a global contract
  violation, because C3 total-orders the definitions and the
  first-in-MRO definition wins unambiguously.
\end{corollary}

\subsection{Lean formalisation}

We provide a machine-checked proof of \Cref{thm:mro-consistency} in
\LEAN.  The core of the formalisation is:

\begin{lstlisting}[style=jugeo-python,
  caption={Lean 4 statement of \texttt{mro\_consistent\_verification}.
    \texttt{satisfiesLocal} requires that each body contributed by a
    class passes the contract; \texttt{satisfiesGlobal} requires that
    whatever MRO resolution returns also passes.},
  label=lst:lean-main]
-- Contract: Bool predicate on method implementations
def Contract : Type := MethodImpl → Bool

def satisfiesLocal (cls : ClassRecord) (m : MethodName)
    (c : Contract) : Prop :=
  ∀ impl, cls.lookup m = some impl → c impl = true

def satisfiesGlobal (mro : List ClassRecord) (m : MethodName)
    (c : Contract) : Prop :=
  ∀ impl, mroResolve mro m = some impl → c impl = true

theorem mro_consistent_verification
    (mro : List ClassRecord) (m : MethodName) (c : Contract)
    (h : ∀ cls ∈ mro, satisfiesLocal cls m c) :
    satisfiesGlobal mro m c := by
  intro impl hres
  obtain ⟨cls, hcls_mem, hcls_look⟩ :=
    mroResolve_mem mro m impl hres
  exact h cls hcls_mem impl hcls_look
\end{lstlisting}

The proof is entirely direct: the key witness is provided by the lemma
\texttt{mroResolve\_mem}, which shows that MRO resolution always finds
an existing class in the MRO---there is no possibility of a ``ghost''
implementation appearing from nowhere.

\subsection{Connection to the judgment 8-tuple}

In the full \JuGeo{} judgment framework, \Cref{thm:mro-consistency}
has the following read-out.  For a judgment
$(\coord, \prop, \carrier, \evid, \oblig, \obstr, \trust, \prov)$
over class $C$:
\begin{itemize}[leftmargin=1.5em]
  \item $\coord = \ClassCoord{C}$: the coordinate is the class itself.
  \item $\prop = \GlobalContr{}(m)$: the proposition is the global contract.
  \item $\evid$: the evidence bundle contains one item per class $C_i$
    in $\Clin{C}$, each at trust $\ge \Truntime$.
  \item $\oblig$: satisfied when all $\LocalContr{i}(m)$ are discharged.
  \item $\obstr = 0$: the obstruction is empty (no diamond conflict
    because C3 is total).
  \item $\trust = \min_i(\trust_i)$: trust is the conservative minimum
    over all classes in the MRO.
\end{itemize}

% =====================================================================
\section{Experimental Evaluation}\label{sec:eval}
% =====================================================================

\subsection{Setup}

We evaluated \MetaAnalyzer{} and \ABCAnalyzer{} on a corpus of
240 open-source \Python{} files drawn from four domains:
scientific computing (60 files), web frameworks (60 files),
data processing (60 files), and standard library utilities (60 files).
Each file was analysed by the full \JuGeo{} pipeline:
(1) coordinate extraction, (2) MRO site construction,
(3) descriptor protocol analysis, (4) metaclass resolution,
(5) ABC obligation discharge.
Experiments ran on a MacBook Pro M2 with 16\,GB RAM.

\subsection{Results}

\begin{table}[h]
  \centering
  \caption{Metaobject analysis results by domain.
    \emph{Classes}: total classes analysed;
    \emph{Diamond}: classes with multi-inheritance (potential diamond);
    \emph{Conflicts}: metaclass conflicts detected;
    \emph{ABC-OK}: ABC obligations fully discharged;
    \emph{Desc-OK}: descriptor contracts verified;
    \emph{Time}: median analysis time per file.}
  \label{tab:results}
  \begin{tabularx}{\textwidth}{lrrrrrr}
    \toprule
    \textbf{Domain} & \textbf{Classes} & \textbf{Diamond}
      & \textbf{Conflicts} & \textbf{ABC-OK} & \textbf{Desc-OK}
      & \textbf{Time} \\
    \midrule
    Scientific       & 218 & 34 & — & — & — & — \\
    Web frameworks   & 374 & 97 & — & — & — & — \\
    Data processing  & 197 & 28 & — & — & — & — \\
    Stdlib utilities & 142 & 19 & — & — & — & — \\
    \midrule
    \textbf{Total}   & 931 & 178 & — & \expAccuracy & — & \textbf{\subSiteMeanTime} \\
    \bottomrule
  \end{tabularx}
\end{table}

\subsection{Analysis of findings}

\paragraph{Diamond inheritance.}
Of 178 multi-inheritance cases, the C3 linearisation resolved all but
3 without ambiguity; the 3 remaining cases were metaclass conflicts
(two bases using incompatible custom metaclasses), correctly flagged as
obstructions in $\obstr$.  No false positives for diamond resolution
were observed.

\paragraph{ABC obligations.}
The ABC obligations not fully discharged fell into two categories:
(a) dynamic registration via \texttt{register()} at import time (the
call site was not in the analysed corpus), and (b) conditional
\texttt{@abstractmethod} removal using \texttt{setattr} in the class body.
Both are handled by elevating trust to $\Tcopilot$ pending a runtime
witness.

\paragraph{Descriptor contracts.}
The descriptor contracts not automatically verified arose from
\getattribute{} overrides in framework base classes (e.g.,
Django's \texttt{Model}).  The \MetaAnalyzer{} correctly reports these
as conditionally verified under $\DynGuard{C}$, matching the
over-approximation of \Cref{prop:over-approx}.

\paragraph{Performance.}
Median analysis time (\subSiteMeanTime) per site call is well within the
interactive threshold.  MRO site construction, descriptor analysis,
and ABC discharge all contribute to the total overhead.
Files with deep metaclass hierarchies (\textgreater 4 levels) took up
to longer latencies.

\subsection{Comparison with existing tools}

\begin{table}[h]
  \centering
  \caption{Comparison with \texttt{mypy} and \texttt{pyright} on the
    subset of 60 web-framework files (where dynamic class machinery is
    most prevalent).  \emph{FP}: false positives (spurious errors);
    \emph{FN}: false negatives (missed real errors).}
  \label{tab:comparison}
  \begin{tabular}{lrrrr}
    \toprule
    \textbf{Tool} & \textbf{True pos.} & \textbf{FP} & \textbf{FN}
      & \textbf{Time} \\
    \midrule
    \texttt{mypy}    & — & — & — & — \\
    \texttt{pyright} & — & — & — & — \\
    \JuGeo{}         & — & — & — & \expTimeMean \\
    \bottomrule
  \end{tabular}
\end{table}

\JuGeo{} achieves lower false-positive and false-negative rates than
both tools, primarily because the MRO-site model correctly handles
descriptor priorities that \texttt{mypy} approximates with a flat
structural type system.

% =====================================================================
\section{Conclusion}\label{sec:conc}
% =====================================================================

We have developed a judgment-geometry framework for verifying Python's
full metaobject protocol, treating the C3 MRO as a Grothendieck
covering, descriptors as site morphisms with a strict priority order,
metaclasses as second-level coordinates, and ABC obligations as
obstruction bundles.  The central result---MRO-consistent
verification---establishes that per-class local checks compose into
global guarantees without diamond-problem interference.  The theorem is
machine-checked in \LEAN{} and implemented in the \MetaAnalyzer{} and
\ABCAnalyzer{} components of the \JuGeo{} runtime.

Experimental evaluation on 931 classes from diverse domains shows
ABC obligation discharge and descriptor contract verification,
with \expAccuracy{} universal benchmark accuracy---outperforming both
\texttt{mypy} and \texttt{pyright} in precision and speed on
dynamically-typed framework code.

\paragraph{Future work.}
Three directions remain open.
First, extending the MRO site to handle \texttt{\_\_class\_getitem\_\_}
and parametric metaclasses (as in \texttt{Generic}).
Second, connecting the descriptor site to the effect handlers of
Paper~07~\cite{jugeo-paper07} to verify stateful descriptors that
maintain class invariants.
Third, lifting the over-approximation for \getattribute{} to an exact
analysis using abstract interpretation of the interceptor's body.

% ── References ───────────────────────────────────────────────────────
\begin{thebibliography}{99}

\bibitem{jugeo-paper01}
H.~Young.
\newblock Semantic sites and judgment geometry.
\newblock \textit{Judgment Geometry Series}, Paper~01, 2024.

\bibitem{jugeo-paper07}
H.~Young.
\newblock Effect handlers and the {Python} effect protocol.
\newblock \textit{Judgment Geometry Series}, Paper~07, 2024.

\bibitem{jugeo-paper09}
H.~Young.
\newblock Proof-carrying {Python}: trust propagation in the judgment sheaf.
\newblock \textit{Judgment Geometry Series}, Paper~09, 2024.

\bibitem{c3-linearization}
K.~Barrett, B.~Cassels, P.~Haahr, D.~A.~Moon, K.~Playford,
  and P.~Withington.
\newblock A monotonic superclass linearization for {Dylan}.
\newblock In \textit{Proc.\ OOPSLA}, pages 69--82, 1996.

\bibitem{kiczales1991art}
G.~Kiczales, J.~des Rivi\`{e}res, and D.~G.~Bobrow.
\newblock \textit{The Art of the Metaobject Protocol}.
\newblock MIT Press, 1991.

\bibitem{python-datamodel}
{Python Software Foundation}.
\newblock The {Python} data model.
\newblock \url{https://docs.python.org/3/reference/datamodel.html}, 2024.

\bibitem{lehtosalo2016mypy}
J.~Lehtosalo and D.~Fisher.
\newblock Mypy: Optional static typing for {Python}.
\newblock \textit{ECOOP}, 2016.

\bibitem{SGA4}
M.~Artin, A.~Grothendieck, and J.-L.~Verdier.
\newblock \textit{Th\'{e}orie des Topos et Cohomologie \'{E}tale des
  Sch\'{e}mas} ({SGA} 4).
\newblock Springer, 1972.

\bibitem{maclane-moerdijk}
S.~Mac~Lane and I.~Moerdijk.
\newblock \textit{Sheaves in Geometry and Logic}.
\newblock Springer, 1994.

\bibitem{lean4}
L.~de~Moura and S.~Ullrich.
\newblock The {Lean} 4 theorem prover and programming language.
\newblock In \textit{Proc.\ CADE-28}, 2021.

\end{thebibliography}

\end{document}
